\subsection{Attention Economics: The Scarcity of Awareness}
\label{subsec:11-4-attention-economics}

%%%%%%%%%%%%%%%%%%%%%%%%%%%%%%%%%%%%%%%%%%%%%%%%%%%%%%%%%%%%%%%%%%%%%%%%%%%%%%%
% SECTION 11.04: Attention Economics
% Position: After Motivated Beliefs (11.03), Before System Levels (11.05)
% Content: Rational Inattention (Sims), Sparse Attention (Gabaix), 
%          Visceral Factors (Loewenstein)
% Rationale: Explains the DETERMINATION of A^base before modulation
%%%%%%%%%%%%%%%%%%%%%%%%%%%%%%%%%%%%%%%%%%%%%%%%%%%%%%%%%%%%%%%%%%%%%%%%%%%%%%%

Section~\ref{subsec:11-3-motivated-beliefs} established how beliefs filter 
awareness through the Belief Filter $B$. But this presupposes that something 
enters base awareness $A^{base}$ in the first place. This section addresses 
the prior question: \emph{What determines the allocation of attention?}

The answer draws on three complementary research programs:
\begin{enumerate}
\item \textbf{Rational Inattention} (Sims): Attention as information-theoretic constraint
\item \textbf{Sparse Attention} (Gabaix): Attention as bounded rationality
\item \textbf{Visceral Factors} (Loewenstein): Attention capture by drive states
\end{enumerate}

\begin{tcolorbox}[colback=red!5!white, colframe=red!75!black, title=Why This Section Matters]
The Awareness function $A$ is not exogenous. People \emph{choose}---implicitly 
or explicitly---what to attend to. Understanding attention allocation is 
prerequisite to understanding awareness at the Moment of Truth.
\end{tcolorbox}

%------------------------------------------------------------------------------
\subsubsection{Rational Inattention: The Sims Framework}
\label{subsubsec:rational-inattention}
%------------------------------------------------------------------------------

\paragraph{The Core Idea.}
Christopher Sims (2003) introduced \textbf{rational inattention} into economics: 
agents have finite information-processing capacity, measured in bits per second 
(Shannon entropy). Attention is scarce and must be allocated optimally.

\begin{tcolorbox}[colback=blue!5!white, colframe=blue!75!black, title=Sims (2003): Implications of Rational Inattention]
``A constraint that actions can depend on observations only through a 
communication channel with finite Shannon capacity is shown to be able to 
play a role similar to that of adjustment costs.''
\end{tcolorbox}

\paragraph{The Information-Theoretic Setup.}
Let $X$ be the state of the world and $S$ the signal the agent processes. 
The information flow is bounded by channel capacity $\kappa$:

\begin{equation}
\boxed{I(X; S) \leq \kappa}
\label{eq:capacity-constraint}
\end{equation}

where $I(X; S)$ is the mutual information between state and signal:

\begin{equation}
I(X; S) = H(X) - H(X|S)
\end{equation}

\begin{itemize}
\item $H(X)$ = entropy of prior beliefs about the state
\item $H(X|S)$ = residual uncertainty after observing signal
\item $\kappa$ = channel capacity (bits per unit time)
\end{itemize}

\paragraph{Implications for Awareness.}
Under rational inattention:

\begin{enumerate}
\item \textbf{Selective attention}: Agents attend to high-value information first
\item \textbf{Coarse perception}: States are perceived with error proportional to $1/\kappa$
\item \textbf{Endogenous inertia}: Behavior is sluggish because tracking requires capacity
\item \textbf{State-dependent precision}: More attention allocated where stakes are higher
\end{enumerate}

\paragraph{Optimal Attention Allocation.}
The agent solves:

\begin{equation}
\max_{f(S|X)} \mathbb{E}[U(a(S), X)] \quad \text{s.t.} \quad I(X;S) \leq \kappa
\label{eq:ri-optimization}
\end{equation}

where $a(S)$ is the action conditional on signal $S$.

\paragraph{Key Result: Logit Choice.}
Matějka \& McKay (2015) show that under rational inattention with discrete 
choices, optimal behavior takes a \textbf{logit form}:

\begin{equation}
Pr(a = a_j) = \frac{\exp(u_j / \lambda)}{\sum_k \exp(u_k / \lambda)}
\end{equation}

where $\lambda$ is inversely related to attention capacity. This provides a 
microfoundation for random utility models.

%------------------------------------------------------------------------------
\subsubsection{Connection to EBF Awareness}
\label{subsubsec:ri-bcm-connection}
%------------------------------------------------------------------------------

Rational inattention determines $A^{base}$:

\begin{equation}
A^{base}_X = f\left(\frac{\partial \mathbb{E}[U]}{\partial I_X}, \kappa\right)
\end{equation}

\paragraph{Interpretation.}
Base awareness of utility component $X$ depends on:
\begin{itemize}
\item The marginal value of information about $X$
\item The overall capacity constraint $\kappa$
\item Competition from other components for limited attention
\end{itemize}

\begin{table}[htbp]
\centering
\caption{Rational Inattention and Awareness Components}
\label{tab:ri-awareness}
\renewcommand{\arraystretch}{1.3}
\begin{tabular}{@{}lll@{}}
\toprule
\textbf{Component} & \textbf{Attention Allocation} & \textbf{Implication for $A^{base}$} \\
\midrule
$U^{INU}_{t=0}$ (instant) & High (immediate feedback) & $A^{base}$ high \\
$U^{INU}_{t=3,4}$ (long-term) & Low (distant, uncertain) & $A^{base}$ low \\
$U^{KNU}$ (collective) & Low (diffuse stakes) & $A^{base}$ low \\
$U^{IDN}$ (identity) & Variable (depends on threat) & $A^{base}$ context-dependent \\
\bottomrule
\end{tabular}
\end{table}

%------------------------------------------------------------------------------
\subsubsection{Sparse Attention: The Gabaix Framework}
\label{subsubsec:sparse-attention}
%------------------------------------------------------------------------------

\paragraph{Sparsity-Based Bounded Rationality.}
Xavier Gabaix (2014, 2019) develops a complementary approach: agents simplify 
their mental models by setting some parameters to zero (``sparsity'').

\begin{tcolorbox}[colback=blue!5!white, colframe=blue!75!black, title=Gabaix (2014): A Sparsity-Based Model of Bounded Rationality]
``The agent pays attention to a variable if and only if the stakes are 
high enough relative to the cost of attention.''
\end{tcolorbox}

\paragraph{The Sparse Max Operator.}
Instead of maximizing $U(a, \mathbf{x})$ over all state variables $\mathbf{x}$, 
the agent uses a sparse representation $\mathbf{x}^s$:

\begin{equation}
x^s_i = 
\begin{cases}
x_i & \text{if } |x_i - x^d_i| \cdot \left|\frac{\partial^2 U}{\partial a \partial x_i}\right| > m_i \\
x^d_i & \text{otherwise}
\end{cases}
\end{equation}

where:
\begin{itemize}
\item $x^d_i$ = default/anchor value
\item $m_i$ = attention cost for variable $i$
\item The condition checks if attending to $x_i$ is ``worth it''
\end{itemize}

\paragraph{Implications.}
\begin{enumerate}
\item \textbf{Anchoring}: Unattended variables stay at defaults
\item \textbf{Threshold effects}: Small deviations are ignored
\item \textbf{Stakes-dependent attention}: High-stakes variables get attention
\item \textbf{Context effects}: Defaults matter when attention is scarce
\end{enumerate}

\paragraph{Integration with Awareness.}
Gabaix's sparsity maps directly to awareness:

\begin{equation}
A^{base}_i = \mathbb{1}\left[|x_i - x^d_i| \cdot \left|\frac{\partial^2 U}{\partial a \partial x_i}\right| > m_i\right]
\end{equation}

This is a \textbf{binary} attention model (attend or not), in contrast to 
Sims's continuous precision model.

%------------------------------------------------------------------------------
\subsubsection{Visceral Factors: The Loewenstein Framework}
\label{subsubsec:visceral-factors}
%------------------------------------------------------------------------------

\paragraph{Attention Capture by Drive States.}
George Loewenstein's research program emphasizes that attention is not always 
under voluntary control. \textbf{Visceral factors}---hunger, thirst, pain, 
fear, arousal---can \emph{capture} attention involuntarily.

\paragraph{Key Papers.}

\begin{table}[htbp]
\centering
\caption{Loewenstein's Contributions to Attention Economics}
\label{tab:loewenstein-papers}
\renewcommand{\arraystretch}{1.3}
\small
\begin{tabular}{@{}llp{5cm}@{}}
\toprule
\textbf{Year} & \textbf{Paper} & \textbf{Core Contribution} \\
\midrule
2000 & Emotions in Economic Theory & Visceral factors as attention drivers \\
2001 & Risk as Feelings (with Weber) & Emotional vs. cognitive risk processing \\
2003 & Time and Decision (with Read) & Attention and intertemporal choice \\
2005 & Hot-Cold Empathy Gap (with Ariely) & State-dependent attention \\
2007 & Exotic Preferences & Projection bias, adaptation \\
\bottomrule
\end{tabular}
\end{table}

\paragraph{The Visceral Factor Model.}
Loewenstein (2000) identifies characteristics of visceral factors:

\begin{enumerate}
\item \textbf{Direct hedonic impact}: They feel good or bad immediately
\item \textbf{Attention capture}: They dominate working memory
\item \textbf{Behavioral override}: They can override deliberative preferences
\item \textbf{State-dependence}: Their influence varies with intensity
\item \textbf{Underestimation}: In ``cold'' states, their future influence is underestimated
\end{enumerate}

\paragraph{Formalization.}
Let $v \in [0,1]$ be visceral intensity. Attention is captured proportionally:

\begin{equation}
A^{visceral} = g(v) \quad \text{with } g'(v) > 0, \; g(0) = 0, \; g(1) = 1
\end{equation}

When $A^{visceral}$ is high, other awareness components are crowded out:

\begin{equation}
A^{base}_{other} = (1 - A^{visceral}) \cdot A^{base,0}_{other}
\end{equation}

\paragraph{The Hot-Cold Empathy Gap.}
A crucial finding (Loewenstein \& Ariely, 2005): People in ``cold'' states 
systematically underestimate how visceral factors will affect their future 
behavior. This creates projection bias in attention:

\begin{equation}
\hat{A}^{visceral}_{future} = A^{visceral}_{current} + (1-\alpha)(A^{visceral}_{true,future} - A^{visceral}_{current})
\end{equation}

where $\alpha < 1$ is the projection bias parameter.

%------------------------------------------------------------------------------
\subsubsection{Integration: Three Mechanisms of Attention Allocation}
\label{subsubsec:integration}
%------------------------------------------------------------------------------

The three frameworks are complementary:

\begin{table}[htbp]
\centering
\caption{Three Mechanisms of Attention Allocation}
\label{tab:three-mechanisms}
\renewcommand{\arraystretch}{1.4}
\begin{tabular}{@{}llll@{}}
\toprule
\textbf{Mechanism} & \textbf{Framework} & \textbf{Nature} & \textbf{Control} \\
\midrule
Optimal allocation & Rational Inattention & Continuous precision & Deliberate \\
Threshold-based & Sparse Attention & Binary (attend/not) & Semi-deliberate \\
Capture & Visceral Factors & Involuntary & Automatic \\
\bottomrule
\end{tabular}
\end{table}

\paragraph{The Complete Attention Model.}
Combining all three:

\begin{equation}
\boxed{
A^{base}_X = \underbrace{(1 - A^{visceral})}_{\text{Available capacity}} \times 
\underbrace{\mathbb{1}[stakes_X > m_X]}_{\text{Sparsity threshold}} \times 
\underbrace{\pi_X(\kappa)}_{\text{RI precision}}
}
\label{eq:complete-attention}
\end{equation}

where:
\begin{itemize}
\item $A^{visceral}$ = attention captured by visceral states
\item $\mathbb{1}[\cdot]$ = sparsity indicator (Gabaix)
\item $\pi_X(\kappa)$ = precision allocated to $X$ under capacity $\kappa$ (Sims)
\end{itemize}

%------------------------------------------------------------------------------
\subsubsection{Attention and the EBF System Levels}
\label{subsubsec:attention-systems}
%------------------------------------------------------------------------------

The attention frameworks map onto EBF's system levels 
(Section~\ref{subsec:11-5-system-levels}):

\begin{table}[htbp]
\centering
\caption{Attention Mechanisms by System Level}
\label{tab:attention-systems}
\renewcommand{\arraystretch}{1.4}
\begin{tabular}{@{}llll@{}}
\toprule
\textbf{System} & \textbf{Attention Type} & \textbf{Framework} & \textbf{Characteristic} \\
\midrule
SYS 0 & Visceral capture & Loewenstein & Automatic, overwhelming \\
SYS 1 & Heuristic allocation & Gabaix (sparse) & Fast, threshold-based \\
SYS 2 & Optimal allocation & Sims (RI) & Slow, deliberate \\
\bottomrule
\end{tabular}
\end{table}

\paragraph{The Hyper-Dissonance Revisited.}
When SYS 0 captures attention ($A^{visceral} \rightarrow 1$), capacity for 
SYS 2 processing collapses. This is the attention-theoretic foundation 
for the Hyper-Dissonance mechanism introduced in 
Section~\ref{subsec:11-5-system-levels}.

%------------------------------------------------------------------------------
\subsubsection{Empirical Evidence}
\label{subsubsec:empirical-evidence}
%------------------------------------------------------------------------------

\paragraph{Evidence for Rational Inattention.}
\begin{itemize}
\item Price stickiness (Maćkowiak \& Wiederholt, 2009): Firms allocate more 
      attention to idiosyncratic shocks than aggregate shocks
\item Portfolio choice (Van Nieuwerburgh \& Veldkamp, 2010): Investors 
      overweight home assets due to attention costs
\item Consumption sluggishness (Luo, 2008): Permanent income behavior with 
      delayed response to income shocks
\end{itemize}

\paragraph{Evidence for Sparse Attention.}
\begin{itemize}
\item Anchoring effects (Tversky \& Kahneman, 1974): Unattended information 
      stays at default
\item Inattention to taxes (Chetty et al., 2009): Small taxes are ignored
\item Shrouded attributes (Gabaix \& Laibson, 2006): Hidden fees not processed
\end{itemize}

\paragraph{Evidence for Visceral Capture.}
\begin{itemize}
\item Sexual arousal effects (Ariely \& Loewenstein, 2006): Risk-taking 
      increases in aroused state
\item Pain and decision-making (Eccleston \& Crombez, 1999): Pain captures 
      attention, degrades decisions
\item Hunger effects (Danziger et al., 2011): Judges more lenient after meals
\end{itemize}

%------------------------------------------------------------------------------
\subsubsection{Implications for Behavioral Intervention}
\label{subsubsec:intervention-implications}
%------------------------------------------------------------------------------

\begin{table}[htbp]
\centering
\caption{Intervention Strategies by Attention Mechanism}
\label{tab:intervention-attention}
\renewcommand{\arraystretch}{1.4}
\begin{tabular}{@{}lp{5cm}p{5cm}@{}}
\toprule
\textbf{Mechanism} & \textbf{Barrier} & \textbf{Intervention Strategy} \\
\midrule
Rational Inattention & 
Limited capacity spread thin &
Reduce complexity; make key info salient \\
Sparse Attention & 
Below threshold, ignored &
Increase stakes visibility; change defaults \\
Visceral Capture & 
Drive states dominate &
Cooling-off periods; pre-commitment \\
\bottomrule
\end{tabular}
\end{table}

\paragraph{The Attention Budget Principle.}
Interventions must respect the attention budget. Adding information can 
\emph{reduce} awareness of key factors by crowding out attention:

\begin{equation}
\frac{\partial A^{base}_{key}}{\partial \text{Info}_{other}} < 0 \quad \text{(attention crowding)}
\end{equation}

This explains why simpler disclosures often outperform detailed ones.

%------------------------------------------------------------------------------
\subsubsection{Axiomatization}
\label{subsubsec:axiomatization}
%------------------------------------------------------------------------------

\begin{tcolorbox}[colback=gray!5!white, colframe=gray!75!black, title=Axiom AWX-A07: Attention Scarcity]
Total attention is bounded by channel capacity:
\begin{equation}
\sum_X A^{base}_X \cdot H(U_X) \leq \kappa
\end{equation}
\end{tcolorbox}

\begin{tcolorbox}[colback=gray!5!white, colframe=gray!75!black, title=Axiom AWX-A08: Stakes-Dependent Attention]
Attention allocation is proportional to decision stakes:
\begin{equation}
A^{base}_X \propto \left|\frac{\partial \mathbb{E}[U]}{\partial I_X}\right|
\end{equation}
\end{tcolorbox}

\begin{tcolorbox}[colback=gray!5!white, colframe=gray!75!black, title=Axiom AWX-A09: Sparsity Threshold]
Below a threshold, information is not processed:
\begin{equation}
A^{base}_X = 0 \quad \text{if } stakes_X < m_X
\end{equation}
\end{tcolorbox}

\begin{tcolorbox}[colback=gray!5!white, colframe=gray!75!black, title=Axiom AWX-A10: Visceral Capture]
Visceral factors crowd out deliberative attention:
\begin{equation}
A^{SYS2} = (1 - A^{visceral}) \cdot A^{SYS2,max}
\end{equation}
\end{tcolorbox}

%------------------------------------------------------------------------------
\subsubsection{Summary}
\label{subsubsec:summary}
%------------------------------------------------------------------------------

This section established that base awareness $A^{base}$ is determined by 
attention allocation, which follows from:

\begin{enumerate}
\item \textbf{Rational Inattention} (Sims): Finite capacity $\kappa$ forces 
      optimal allocation across information sources
\item \textbf{Sparse Attention} (Gabaix): Below-threshold information is 
      ignored entirely (set to default)
\item \textbf{Visceral Capture} (Loewenstein): Drive states can hijack 
      attention involuntarily
\end{enumerate}

The complete model of base awareness is:

\begin{equation}
\boxed{
A^{base}_X = (1 - A^{visceral}) \times \mathbb{1}[stakes_X > m_X] \times \pi_X(\kappa)
}
\end{equation}

This provides the foundation for understanding \emph{what enters awareness} 
before the Belief Filter (Section~\ref{subsec:11-3-motivated-beliefs}) 
determines \emph{what is accepted}.

The full Awareness function is now:

\begin{equation}
\boxed{
A_X = \underbrace{(1 - A^{visceral}) \times \mathbb{1}[stakes > m] \times \pi(\kappa)}_{\text{Attention Economics } (A^{base})} \times \underbrace{B(\hat{\theta}, signal, U_{IDN})}_{\text{Belief Filter}}
}
\end{equation}

%------------------------------------------------------------------------------
\subsubsection{BEATRIX Module Reference}
\label{subsubsec:beatrix-ref}
%------------------------------------------------------------------------------

\begin{table}[htbp]
\centering
\caption{Section 11.04 BEATRIX References}
\label{tab:11-4-beatrix}
\renewcommand{\arraystretch}{1.3}
\begin{tabular}{@{}lll@{}}
\toprule
\textbf{Element} & \textbf{Module} & \textbf{Reference} \\
\midrule
Attention Scarcity & AWX-A07 & Capacity Constraint \\
Stakes-Dependent Attention & AWX-A08 & Optimal Allocation \\
Sparsity Threshold & AWX-A09 & Gabaix Threshold \\
Visceral Capture & AWX-A10 & Loewenstein Override \\
\bottomrule
\end{tabular}
\end{table}

%%%%%%%%%%%%%%%%%%%%%%%%%%%%%%%%%%%%%%%%%%%%%%%%%%%%%%%%%%%%%%%%%%%%%%%%%%%%%%%
% END OF SECTION 11.04
%%%%%%%%%%%%%%%%%%%%%%%%%%%%%%%%%%%%%%%%%%%%%%%%%%%%%%%%%%%%%%%%%%%%%%%%%%%%%%%
